%=========================
% Sections/01-GioiThieuDeTai.tex
% Chương 1 – Giới thiệu đề tài
%=========================

\section{Giới thiệu đề tài}

\subsection{Lý do chọn đề tài}
\indent Trong thương mại điện tử, người dùng thường mô tả nhu cầu bằng ngôn ngữ tự nhiên, hình ảnh tham chiếu hoặc kết hợp cả hai thay vì nhập chính xác tên sản phẩm. Tìm kiếm dựa chủ yếu trên từ khóa khó khai thác đồng thời ý nghĩa truy vấn, đặc trưng hình ảnh và các ràng buộc như danh mục, màu sắc, kiểu dáng hoặc mức giá, nên kết quả có thể chưa phù hợp với ý định mua sắm.

\indent Vì vậy, đề tài \textit{\enquote{Xây dựng Trợ lý Mua sắm Đa phương tiện sử dụng RAG và Computer Vision}} được lựa chọn để hỗ trợ tìm kiếm bằng văn bản, hình ảnh, truy vấn kết hợp và tư vấn qua hội thoại có ngữ cảnh. Việc kết hợp RAG với thị giác máy tính giúp phản hồi dựa trên dữ liệu sản phẩm có thể kiểm chứng, đồng thời tạo cơ sở tích hợp AI với các nghiệp vụ thương mại điện tử mà không thay thế quy tắc xác nhận giá, tồn kho và đơn hàng.

\subsection{Mục tiêu}
\indent Mục tiêu tổng quát của đề tài là xây dựng, triển khai và đánh giá một trợ lý mua sắm đa phương thức tích hợp trong hệ thống thương mại điện tử. Các mục tiêu cụ thể gồm:
\begin{itemize}
    \item Xây dựng các chức năng thương mại điện tử cốt lõi và giao diện phục vụ khách hàng, quản trị viên.
    \item Hỗ trợ tìm kiếm sản phẩm bằng văn bản, hình ảnh và truy vấn kết hợp có ràng buộc.
    \item Ứng dụng RAG và bộ nhớ hội thoại để tư vấn dựa trên dữ liệu sản phẩm được truy xuất.
    \item Tích hợp Frontend, Backend, AI Service, MySQL và Qdrant; triển khai hệ thống trên môi trường đám mây.
    \item Kiểm thử các luồng nghiệp vụ và đánh giá chất lượng truy hồi, phản hồi hội thoại, độ trễ, bảo mật và khả dụng.
\end{itemize}

\subsection{Phạm vi}
\indent Đề tài tập trung vào một hệ thống thương mại điện tử trong lĩnh vực thời trang, được triển khai và đánh giá ở quy mô thử nghiệm trên môi trường cục bộ và Google Cloud.

\begin{itemize}
    \item \textbf{Trong phạm vi:} quản lý sản phẩm và người dùng; giỏ hàng và đặt hàng COD; tìm kiếm bằng văn bản, hình ảnh và kết hợp; tư vấn bằng RAG có bộ nhớ hội thoại; quản trị hệ thống; đồng bộ dữ liệu MySQL--Qdrant và kiểm thử các chức năng chính.
    \item \textbf{Ngoài phạm vi:} thanh toán trực tuyến thực tế, tích hợp đơn vị vận chuyển, vận hành thương mại ở quy mô lớn và ứng dụng di động.
    \item \textbf{Điều kiện thử nghiệm:} dữ liệu sản phẩm đã được chuẩn hóa; các mô hình và dịch vụ bên ngoài sẵn sàng; kết quả được ghi nhận trong phạm vi dữ liệu và hạ tầng của đề tài.
\end{itemize}

\subsection{Cấu trúc báo cáo}
\indent Báo cáo hiện gồm tám chương.\\
Chương 1 giới thiệu lý do chọn đề tài, mục tiêu, phạm vi và cấu trúc báo cáo.\\
Chương 2 phân tích các bên liên quan, quy trình nghiệp vụ và yêu cầu hệ thống.\\
Chương 3 trình bày kiến thức nền tảng.\\
Chương 4 giới thiệu các công nghệ được lựa chọn.\\
Chương 5 trình bày thiết kế kiến trúc, dữ liệu và các chức năng chính.\\
Chương 6 trình bày quá trình hiện thực, cấu hình và triển khai hệ thống.\\
Chương 7 trình bày phương pháp, kết quả kiểm thử và đánh giá hệ thống.\\
Chương 8 tổng kết kết quả đạt được, hạn chế và đề xuất hướng phát triển trong tương lai.
